\section{Platform (secondary) objectives}

The analytical objectives above describe what the system computes; the platform objectives describe what the users can see and do. HospitAI defines three user roles: patient, doctor, and manager, each with its own interface and its own view of the data.

\vspace{0.5em}
\noindent\textbf{Patient Interface}

In HospitAI, a patient can sign in and view only their own medical record. The patient's user interface displays the individual's profile (which includes name, age, marital status, and contact information), a summary of key health parameters (QALY, BMI, conditions, number of medications, number of emergencies, days since the last visit), and a timeline that records all visits, conditions, medications, procedures, immunizations, allergies, and care plans. A patient cannot view any other individual's record, does not have access to any clustering and sequence mining results, and cannot make any queries to the database.

\vspace{0.5em}
\noindent\textbf{Doctor Interface}

The doctor has the ability to browse the patient directory, filter the patients according to their gender, age range, race, marital status, or clinical status, and access the details page for any of the patients. On the patient detail view, the doctor will have all the clinical tabs of the patient, which include conditions, medications, procedures, immunizations, allergies, care plans, and encounters, in addition to two further sections containing the cluster associated with the patient and the sequential patterns associated with the patient. The doctor can also access the cluster summary report that will enable them to understand how the current patient fits within other similar patients.

\vspace{0.5em}
\noindent\textbf{Manager Interface}

Unlike a doctor who only deals with one patient at a time, the manager will be working with aggregate perspectives. The manager will work with aggregate views such as the home dashboard containing KPIs and charts (number of patients, mortality rate, number of conditions, number of care plans, number of procedures, encounters by week and month, demographics, most common conditions), the cluster browser containing all clusters and their LLM-generated descriptions, the sequence browser showing all sequences and their support values, and the text-to-SQL chat through which one can freely explore the database. It should be noted that the manager can drill down from a cluster to the patients in that cluster and from a sequence to the patients that follow a particular sequence, though the essence of the role remains focused on the aggregate perspective rather than individual patient management.

\vspace{0.5em}
\noindent\textbf{Scope and Non-Goals}

To keep the project within the time frame of this thesis, a few things are explicitly out of scope. First of all, HospitAI does not connect to any real hospitals' databases or to any patient data - it operates with the synthetic dataset only. The platform is not designed as a replacement for the human clinician who can provide treatment advice. Any text generated by the LLM is merely a description of the statistical information and does not constitute a medical recommendation based on that description. Clustering and sequence mining are run as offline batch tasks without any real-time computation, which could be too costly for the research objectives pursued within the thesis framework. Finally, there is no fine-grained permission system implemented across the three user roles at the level of a production-quality authorization module, since the research part of this thesis concerns itself with analysis capabilities only.

Given the above project objectives, Chapter~\ref{ch:research} gives an overview of the literature relevant to this thesis and sets the research in context.
